\documentclass[12pt,a4paper]{article}

% Paquetes necesarios
\usepackage[utf8]{inputenc}
\usepackage[T1]{fontenc}
\usepackage[spanish]{babel}
\usepackage{amsmath}
\usepackage{amssymb}
\usepackage{amsthm}
\usepackage{mathtools}
\usepackage{geometry}
\usepackage{hyperref}
\usepackage{bookmark}
\usepackage{graphicx}
\usepackage{booktabs}
\usepackage{array}
\usepackage{enumitem}
\usepackage{fancyhdr}
\usepackage{titlesec}
\usepackage{textcomp}
\usepackage{tikz}
\usetikzlibrary{arrows.meta, positioning, calc, decorations.pathreplacing, shapes.geometric}
\usepackage{pgfplots}
\pgfplotsset{compat=1.18}
\usepgfplotslibrary{fillbetween}
\usepackage{caption}

% Configuración de página
\geometry{margin=2.5cm}
\setlength{\headheight}{15pt}
\pagestyle{fancy}
\fancyhf{}
\fancyhead[L]{Fundamentos de Valoración de Activos}
\fancyhead[R]{Jesus Villota Miranda - CEMFI}
\fancyfoot[C]{\thepage}

% Configuración de hipervínculos
\hypersetup{
    colorlinks=true,
    linkcolor=blue,
    urlcolor=blue,
    citecolor=blue
}

% Definición de entornos matemáticos
\theoremstyle{definition}
\newtheorem{definicion}{Definición}[section]
\newtheorem{teorema}{Teorema}[section]
\newtheorem{proposicion}{Proposición}[section]
\newtheorem{ejemplo}{Ejemplo}[section]
\newtheorem{corolario}{Corolario}[section]
\newtheorem{ejercicio}{Ejercicio}[section]
\newtheorem{solucion}{Solución}[section]

% Comandos personalizados
\newcommand{\R}{\mathbb{R}}
\newcommand{\N}{\mathbb{N}}
\newcommand{\Z}{\mathbb{Z}}
\newcommand{\C}{\mathbb{C}}
\newcommand{\E}{\mathbb{E}}
\newcommand{\V}{\mathbb{V}}
\newcommand{\Cov}{\mathbb{C}\text{ov}}

% Estilos TikZ para visualizaciones financieras
\tikzset{
  timeline/.style={-Stealth, thick, black},
  cashflow/.style={-Stealth, very thick, blue!70},
  discount/.style={-Stealth, dashed, red!70},
  simple/.style={fill=blue!30, draw=blue!70},
  compound/.style={fill=green!30, draw=green!70},
  timepoint/.style={circle, fill, inner sep=2pt},
  payment/.style={-Stealth, thick, orange!80},
  growth/.style={-Stealth, thick, purple!80}
}

% Configuración de títulos
\titleformat{\section}{\Large\bfseries}{\thesection}{1em}{}
\titleformat{\subsection}{\large\bfseries}{\thesubsection}{1em}{}
\titleformat{\subsubsection}{\normalsize\bfseries}{\thesubsubsection}{1em}{}



% Título del documento
\title{\textbf{Ofertas Públicas Iniciales (OPIs)}\\
\large Initial Public Offerings (IPOs)}
\author{Jesus Villota Miranda\\
CEMFI}
\date{\today}

\begin{document}

\maketitle

\tableofcontents
\newpage

\section{Introducción}

Este documento presenta un tratamiento riguroso y autocontenido de las Ofertas Públicas Iniciales (OPIs), también conocidas como Initial Public Offerings (IPOs). Las OPIs representan uno de los mecanismos más importantes mediante los cuales las empresas privadas acceden a los mercados de capitales públicos, transformando su estructura de propiedad y financiación.

El proceso de salida a bolsa implica una serie de decisiones estratégicas complejas que afectan tanto a la empresa emisora como a sus accionistas existentes y potenciales inversores. Este documento se estructura para proporcionar una comprensión completa de los aspectos teóricos, prácticos y empíricos de las OPIs, desde su definición formal hasta las características empíricas que han desafiado las explicaciones tradicionales de los mercados financieros eficientes.

El documento se organiza en cinco secciones principales: (1) la definición formal de OPIs y emisiones de capital, (2) los beneficios y costes asociados con la salida a bolsa, (3) el proceso previo a la OPI y el papel de los intermediarios financieros, (4) los métodos de fijación de precios, y (5) las tres características empíricas desconcertantes de las OPIs que han generado extensa investigación académica. Cada sección incluye definiciones formales, análisis teóricos rigurosos y evidencia empírica relevante.

\section{Definición de Ofertas Públicas Iniciales}

\subsection{Conceptos Fundamentales}

Las emisiones de capital representan una de las formas principales mediante las cuales las empresas obtienen financiación de los mercados de capitales. Es fundamental distinguir entre diferentes tipos de emisiones según el estado previo de la empresa y la naturaleza de las acciones ofrecidas.

\begin{definicion}[Oferta Pública Inicial (OPI)]
Una \textbf{Oferta Pública Inicial} (OPI, o IPO por sus siglas en inglés) se define como la \textbf{primera oferta de acciones de una empresa al público en general}. Una OPI se caracteriza por dos elementos esenciales:
\begin{enumerate}
\item \textbf{Primera}: La empresa era previamente una empresa privada.
\item \textbf{Pública}: La oferta no es una colocación privada (\textit{private placement}), sino que está dirigida al público en general.
\end{enumerate}
\end{definicion}

La definición de OPI enfatiza que se trata de un evento único en la vida de la empresa: la transición de ser una empresa privada a convertirse en una empresa pública cuyas acciones se negocian en mercados organizados.

\begin{definicion}[Oferta Secundaria de Capital (SEO)]
Una \textbf{Oferta Secundaria de Capital} (SEO, por sus siglas en inglés) es una emisión de acciones por parte de una empresa que \textbf{ya es pública}. Al igual que en una OPI, las acciones emitidas en una SEO pueden ser acciones primarias o secundarias.
\end{definicion}

La distinción entre OPI y SEO es fundamental, ya que las empresas públicas tienen acceso continuo a los mercados de capitales, mientras que las OPIs representan un evento único de transición.

\subsection{Emisiones Primarias y Secundarias}

Dentro del contexto de cualquier emisión de capital (tanto OPI como SEO), es crucial distinguir entre emisiones primarias y secundarias, ya que estas determinan quién recibe los fondos obtenidos y cómo se afecta la estructura de propiedad de la empresa.

\begin{definicion}[Emisión Primaria]
En una \textbf{emisión primaria}, las acciones ofrecidas son acciones nuevas emitidas por la empresa. Los fondos obtenidos de la venta de estas acciones se invierten en la empresa, aumentando el capital social y proporcionando recursos para financiar proyectos de inversión, reducción de deuda, o necesidades operativas.
\end{definicion}

\begin{definicion}[Emisión Secundaria]
En una \textbf{emisión secundaria}, las acciones son vendidas por accionistas actuales de la empresa (por ejemplo, fundadores, capitalistas de riesgo, o inversores institucionales). Los fondos obtenidos de la venta de estas acciones van directamente a los vendedores anteriores, no a la empresa. La empresa no recibe ningún capital adicional de una emisión secundaria pura.
\end{definicion}

En la práctica, muchas OPIs combinan emisiones primarias y secundarias: la empresa emite nuevas acciones para obtener capital, mientras que algunos accionistas existentes venden parte de su participación para realizar ganancias o diversificar sus carteras.

\section{Beneficios y Costes de las OPIs}

La decisión de realizar una OPI implica una evaluación cuidadosa de los beneficios y costes asociados. Esta evaluación es fundamental para determinar si la salida a bolsa maximiza el valor para los accionistas existentes y si es el momento óptimo para realizar la transición.

\subsection{Beneficios de las OPIs}

Los beneficios de realizar una OPI pueden clasificarse en varias categorías: acceso a financiación, mejoras en liquidez y valoración, acceso a información del mercado, y consideraciones estratégicas.

\subsubsection{Acceso a Nuevas Fuentes de Financiación}

\begin{proposicion}[Acceso Ampliado a Financiación]
Una OPI proporciona acceso a nuevas fuentes de financiación cuando las fuentes actuales se han agotado o son insuficientes. Las OPIs pueden servir como sustituto de financiación bancaria, capital de riesgo y business angels.
\end{proposicion}

Para empresas en crecimiento que han agotado las fuentes de financiación privada tradicionales, las OPIs representan una vía para acceder a los mercados públicos de capital, que típicamente ofrecen mayor capacidad de financiación y menores restricciones que las fuentes privadas.

\subsubsection{Aumento de Liquidez y Valoración}

\begin{proposicion}[Efecto de Liquidez sobre el Valor]
El aumento de la liquidez en las acciones de la empresa aumenta su valor. La liquidez reduce el coste de capital al disminuir la prima de iliquidez requerida por los inversores.
\end{proposicion}

Las acciones de empresas públicas son más líquidas que las de empresas privadas, ya que pueden negociarse fácilmente en mercados secundarios. Esta mayor liquidez reduce el coste de capital de la empresa y puede aumentar el valor de las acciones existentes.

\begin{proposicion}[Acceso a Mercados Públicos de Capital]
El acceso a los mercados públicos de capital es más barato que el capital de riesgo gracias a la mayor liquidez. Los inversores públicos requieren menores primas de riesgo que los inversores de capital de riesgo debido a la capacidad de diversificar y la mayor liquidez de sus inversiones.
\end{proposicion}

\subsubsection{Salida para Financiadores Previos}

Las OPIs proporcionan una vía de salida (\textit{exit}) para financiadores previos, como capitalistas de riesgo y business angels, que típicamente tienen horizontes de inversión limitados. Sin embargo, esta salida está sujeta a restricciones:

\begin{itemize}
\item \textbf{Período de bloqueo (\textit{lock-up period})}: Los accionistas existentes generalmente están sujetos a un período durante el cual no pueden vender sus acciones después de la OPI.
\item \textbf{Restricciones regulatorias}: Las regulaciones del mercado pueden imponer restricciones adicionales sobre la venta de acciones por parte de insiders.
\item \textbf{Restricciones de los bancos de inversión}: Los suscriptores pueden imponer restricciones contractuales para estabilizar el precio de las acciones después de la OPI.
\end{itemize}

\subsubsection{Información del Precio de las Acciones}

\begin{proposicion}[Valor Informativo del Precio de las Acciones]
El precio de las acciones agrega información de analistas y del mercado, proporcionando:
\begin{enumerate}
\item Mejores decisiones de inversión basadas en señales del mercado.
\item Evaluación del desempeño gerencial por parte del mercado.
\end{enumerate}
\end{proposicion}

El precio de mercado de las acciones refleja la información agregada de todos los participantes del mercado, proporcionando retroalimentación valiosa sobre las decisiones estratégicas y el desempeño de la gestión.

\subsubsection{Mayor Visibilidad}

La salida a bolsa aumenta significativamente la visibilidad de la empresa, lo cual es particularmente importante para industrias que requieren relaciones a largo plazo con clientes y proveedores. Por ejemplo, empresas de software que requieren servicios continuos o actualizaciones se benefician de la credibilidad y visibilidad asociadas con el estatus de empresa pública.

\subsubsection{Fusiones y Adquisiciones}

Las empresas públicas tienen ventajas en el contexto de fusiones y adquisiciones:

\begin{enumerate}
\item Es más fácil para otras empresas notar y evaluar una empresa pública para posibles sinergias.
\item Las OPIs y SEOs se utilizan frecuentemente para financiar adquisiciones, proporcionando una moneda de cambio (acciones) para transacciones estratégicas.
\end{enumerate}

\subsection{Costes de las OPIs}

A pesar de los beneficios significativos, las OPIs conllevan costes sustanciales que deben ser cuidadosamente considerados en la decisión de salida a bolsa.

\subsubsection{Costes Directos e Indirectos}

\begin{definicion}[Costes Directos de una OPI]
Los \textbf{costes directos} de una OPI incluyen:
\begin{itemize}
\item Comisión de suscripción (\textit{underwriting fee}): aproximadamente 7\% del valor de la emisión.
\item Honorarios legales y de asesoramiento.
\item Impuestos y otros costes regulatorios.
\end{itemize}
\end{definicion}

\begin{definicion}[Costes Indirectos de una OPI]
Los \textbf{costes indirectos} incluyen el coste de oportunidad del tiempo de la dirección ejecutiva dedicado al proceso de OPI, que puede distraer de las operaciones diarias de la empresa.
\end{definicion}

Los costes directos de las OPIs son significativos y varían según el tamaño de la emisión, la reputación del banco de inversión, y las condiciones del mercado. Para emisiones pequeñas, los costes pueden representar una fracción aún mayor del valor total.

\subsubsection{Dispersión de la Propiedad y Pérdida de Control}

La salida a bolsa típicamente resulta en una dispersión significativa de la propiedad, reduciendo el control de los accionistas fundadores y directivos. Esta dispersión puede:

\begin{itemize}
\item Reducir la capacidad de los fundadores para tomar decisiones estratégicas sin consultar a accionistas externos.
\item Aumentar el riesgo de adquisiciones hostiles.
\item Crear conflictos de agencia entre directivos y accionistas dispersos.
\end{itemize}

\subsubsection{Cumplimiento de Requisitos Regulatorios}

Las empresas públicas deben cumplir con todos los requisitos regulatorios aplicables, que incluyen:

\begin{itemize}
\item Presentación regular de estados financieros auditados.
\item Divulgación de información material sobre la empresa.
\item Cumplimiento de normas de gobierno corporativo.
\item Reuniones anuales de accionistas y comunicación continua con inversores.
\end{itemize}

Estos requisitos imponen costes administrativos y de cumplimiento significativos que no existen para empresas privadas.

\subsubsection{Mayor Divulgación y Escrutinio}

Las empresas públicas están sujetas a un nivel de divulgación y escrutinio mucho mayor que las empresas privadas:

\begin{itemize}
\item Deben divulgar información financiera detallada y estratégica.
\item Están sujetas al escrutinio de analistas, medios de comunicación, y reguladores.
\item La información confidencial puede ser más difícil de proteger.
\end{itemize}

\subsubsection{Dependencia de la Percepción Externa}

\begin{proposicion}[Dependencia del Valor en la Percepción Externa]
El valor de la empresa se vuelve más dependiente de la percepción externa. El precio de las acciones refleja las expectativas del mercado, que pueden no estar alineadas con la realidad operativa de la empresa a corto plazo.
\end{proposicion}

Esta dependencia puede crear presión para gestionar las expectativas del mercado y puede llevar a decisiones estratégicas que priorizan resultados a corto plazo sobre valor a largo plazo.

\section{Financiación en Etapas Tempranas}

Antes de considerar una OPI, las empresas típicamente pasan por varias etapas de financiación privada. Comprender estas etapas es esencial para entender el contexto completo del proceso de salida a bolsa.

\subsection{Tipos de Financiación Privada}

\begin{definicion}[Business Angels]
Los \textbf{Business Angels} son inversores privados que invierten su capital personal en empresas en etapas muy tempranas. Proporcionan financiación, experiencia empresarial, y redes de contactos a startups y empresas en fase inicial.
\end{definicion}

\begin{definicion}[Capital de Riesgo (Venture Capital)]
El \textbf{Capital de Riesgo} (VC) invierte en empresas en etapas tempranas o de alto crecimiento, típicamente adquiriendo una participación minoritaria (10-30\%) en la empresa. Los fondos de capital de riesgo proporcionan financiación escalonada en múltiples rondas (Series A, B, C, etc.) a medida que la empresa crece.
\end{definicion}

\begin{definicion}[Deuda Privada]
La \textbf{deuda privada} es una forma de financiación mediante la cual una empresa emite instrumentos de deuda privados. Típicamente es más adecuada para empresas con flujos de caja más estables. A diferencia del capital, la deuda no otorga derechos de control y representa valores de renta fija.
\end{definicion}

Los instrumentos híbridos, como la deuda convertible, combinan características de deuda y capital, permitiendo a los inversores convertir su deuda en acciones bajo ciertas condiciones.

\begin{definicion}[Capital Privado (Private Equity)]
El \textbf{Capital Privado} (PE) no es para etapas tempranas. Invierte en empresas maduras y establecidas que ya están generando ingresos o beneficios. Típicamente adquiere una participación mayoritaria o la propiedad completa (a menudo 100\%) de la empresa.
\end{definicion}

La distinción entre capital de riesgo y capital privado es fundamental: el primero se enfoca en empresas en crecimiento temprano, mientras que el segundo se enfoca en empresas establecidas que requieren reestructuración o crecimiento acelerado.

\section{El Proceso Previo a las OPIs}

El proceso de realizar una OPI es complejo y requiere una preparación cuidadosa que puede extenderse durante varios meses. Este proceso involucra múltiples participantes, incluyendo la empresa emisora, bancos de inversión, abogados, auditores, y reguladores.

\subsection{Linea de Tiempo del Proceso}

El proceso de OPI típicamente sigue una línea de tiempo estructurada que comienza varios meses antes de la fecha planificada de salida a bolsa.

\subsubsection{Etapa Inicial}

En la etapa inicial, la empresa discute la posibilidad de realizar una OPI con bancos de inversión. Durante estas discusiones, la empresa debe considerar:

\begin{itemize}
\item Qué podría valer la empresa en el mercado público.
\item Existe un umbral mínimo de valoración por debajo del cual una OPI puede no ser viable debido a los costes fijos del proceso.
\end{itemize}

\subsubsection{Seis Meses Antes de la OPI Planificada}

Aproximadamente seis meses antes de la fecha planificada de la OPI, la empresa debe:

\begin{enumerate}
\item \textbf{Asegurar que la estructura corporativa sea adecuada}: La estructura legal y organizacional de la empresa debe ser compatible con los requisitos de una empresa pública.
\item \textbf{Preparar estados financieros auditados}: La empresa debe tener estados financieros auditados preparados. Esto puede requerir cambiar de auditor si el auditor actual no cumple con los requisitos para empresas públicas.
\end{enumerate}

\subsubsection{Dos a Tres Meses Antes de la OPI Planificada}

En los dos a tres meses previos a la OPI, el proceso se intensifica:

\begin{enumerate}
\item \textbf{Seleccionar un suscriptor principal (\textit{lead underwriter})}: El banco de inversión principal coordinará el proceso de OPI y será responsable de la suscripción de las acciones.
\item \textbf{Seleccionar un sindicato de analistas y corredores}: El sindicato ayudará en la comercialización de las acciones, preparará informes de investigación, y realizará conferencias con inversores.
\item \textbf{Período de educación de inversores}: Los analistas y la dirección educan a inversores potenciales sobre la empresa, su modelo de negocio, y sus perspectivas futuras.
\item \textbf{Road show}: La dirección ejecutiva realiza presentaciones a inversores institucionales en múltiples ciudades para generar interés en la oferta.
\item \textbf{El prospecto}: Se prepara y distribuye el documento de registro que contiene toda la información relevante sobre la empresa y la oferta.
\item \textbf{Book-building}: El proceso mediante el cual el banco de inversión recopila indicaciones de interés de inversores para determinar la demanda y establecer el precio final.
\end{enumerate}

\subsection{El Prospecto}

El prospecto es el documento legal fundamental en el proceso de OPI. Contiene información detallada que los inversores necesitan para tomar una decisión informada sobre la inversión.

\begin{definicion}[Prospecto]
El \textbf{prospecto} es el documento de registro oficial que contiene todos los detalles de la OPI, incluyendo:
\begin{itemize}
\item Número de acciones ofrecidas.
\item Quién ofrece las acciones (empresa o accionistas existentes).
\item Estructura de propiedad mayoritaria antes y después de la OPI.
\item Discusión del negocio y datos financieros históricos.
\item Informe del contador (auditor).
\item Uso de los fondos obtenidos (\textit{use of proceeds}).
\item Factores de riesgo, incluyendo litigios pendientes.
\item Discusión de la dirección ejecutiva.
\item Discusión de la posibilidad de opción de sobredemanda (\textit{overallotment option} o opción \textit{green-shoe}).
\end{itemize}
\end{definicion}

En los Estados Unidos, los prospectos aparecen en EDGAR, la base de datos de la SEC (Securities and Exchange Commission), proporcionando acceso público a esta información.

\textbf{Nota importante}: Hasta este punto en el proceso, un precio final no ha sido determinado, aunque la empresa y el banco obviamente tienen un precio en mente basado en la valoración preliminar.

\subsection{Contratos con Bancos de Inversión}

Las empresas pueden estructurar sus contratos con bancos de inversión de dos formas principales, cada una con implicaciones diferentes para la distribución del riesgo.

\begin{definicion}[Compromiso Firme (\textit{Firm Commitment})]
En un acuerdo de \textbf{compromiso firme}, el suscriptor compra las acciones y las revende al público. Por tanto, el suscriptor asume la incertidumbre de la emisión. Si las acciones no se venden al precio ofrecido, el suscriptor mantiene el inventario y asume las pérdidas.
\end{definicion}

El método de compromiso firme es el método usual utilizado por bancos de reputación en los Estados Unidos. Bajo este método, el banco garantiza la venta de las acciones al precio acordado, asumiendo el riesgo de mercado.

\begin{definicion}[Mejores Esfuerzos (\textit{Best Efforts})]
En un acuerdo de \textbf{mejores esfuerzos}, el suscriptor no garantiza que las acciones se venderán. Por tanto, la empresa asume la incertidumbre de la emisión. Si la demanda es insuficiente, la oferta puede ser cancelada o reducida.
\end{definicion}

\subsubsection{Compensación del Banco de Inversión}

La estructura de compensación difiere según el tipo de contrato:

\begin{itemize}
\item \textbf{Mejores esfuerzos}: Honorario fijo, independiente de las acciones vendidas. El acuerdo puede ser cancelado si la demanda es insuficiente.
\item \textbf{Compromiso firme}: La ganancia del suscriptor se basa en cuántas acciones vende y en el diferencial (\textit{spread}) entre el precio de compra con descuento y el precio de venta al público.
\end{itemize}

Los diferenciales típicos varían por región:
\begin{itemize}
\item 3.09\% en Asia Pacífico.
\item 3.59\% en Europa.
\item 7\% en América del Norte.
\item 5\% en el Reino Unido.
\end{itemize}

\subsection{Papel del Suscriptor}

El suscriptor principal (banco de inversión) desempeña múltiples funciones críticas en el proceso de OPI.

\begin{definicion}[Funciones del Suscriptor Principal]
El suscriptor principal es un banco de inversión que realiza las siguientes tareas:
\begin{enumerate}
\item Proporciona un monitoreo "independiente" de la calidad de la empresa al mercado, certificando la viabilidad y atractivo de la oferta.
\item Asesora a la empresa en:
   \begin{itemize}
   \item Valoración de la empresa.
   \item Fijación del precio de la emisión.
   \end{itemize}
\item Maneja parte del proceso administrativo, incluyendo la preparación de documentación regulatoria y la coordinación con otros participantes.
\end{enumerate}
\end{definicion}

\subsubsection{Rol Continuo del Suscriptor}

El papel del suscriptor no necesariamente termina después de la emisión de capital:

\begin{enumerate}
\item \textbf{Cobertura de investigación}: Proporcionan cobertura de investigación (informes de analistas) después de la OPI. En el caso de SEOs, generalmente ya proporcionaban cobertura antes.
\item \textbf{Creador de mercado}: El suscriptor principal a menudo actúa como "creador de mercado" (\textit{market-maker}), proporcionando liquidez, especialmente en los primeros días de negociación después de una OPI.
\item \textbf{Estabilización de precios}: En el caso de OPIs, puede ser responsable de "estabilizar" los precios si comienzan a caer por debajo del precio de oferta.
\end{enumerate}

\subsection{Opción de Sobredemanda (Green Shoe)}

La opción de sobredemanda, también conocida como opción \textit{green-shoe}, es un mecanismo importante que permite al suscriptor gestionar el riesgo de estabilización de precios.

\begin{definicion}[Opción de Sobredemanda (\textit{Green Shoe Option})]
La \textbf{opción de sobredemanda} (también conocida como opción \textit{green-shoe}) es una cláusula contractual en el contrato de suscripción que permite al banco de inversión vender más acciones en la OPI de las que realmente se están emitiendo. Dentro de (generalmente) 30 días después de la OPI, el suscriptor tiene la opción de comprar acciones adicionales de la empresa (generalmente un 15\% adicional del tamaño de la OPI) al precio de oferta.
\end{definicion}

\subsubsection{Mecanismo de Funcionamiento}

Con esta opción disponible, el suscriptor a menudo vende realmente el 115\% del tamaño de la emisión en la OPI, dejándose con una posición corta en las acciones:

\begin{itemize}
\item \textbf{Si el precio de las acciones sube}: El suscriptor ejerce la opción, cubriendo así su posición corta a un precio fijo, y el tamaño de la OPI será del 115\%.
\item \textbf{Si el precio de las acciones baja}: El suscriptor cubre la posición corta comprando las acciones excedentes (15\%) en el mercado y no ejerce la opción. La cantidad de acciones en el mercado será del 100\% del tamaño de la OPI.
\end{itemize}

De cualquier manera, el suscriptor evita quedarse con un inventario positivo de las acciones, gestionando así el riesgo de mantener acciones no vendidas.

\section{Métodos de Fijación de Precios}

La determinación del precio de oferta en una OPI es uno de los aspectos más críticos del proceso. Existen tres métodos principales de fijación de precios, cada uno con características distintas y utilizados en diferentes contextos geográficos y de mercado.

\subsection{Tres Métodos de Fijación de Precios}

\begin{definicion}[Ofertas de Precio Fijo]
En las \textbf{ofertas de precio fijo}, el banco de inversión y la empresa determinan conjuntamente el precio al cual se venderán los valores. Si hay exceso de demanda, las acciones típicamente se racionan de forma proporcional (\textit{pro-rata}).
\end{definicion}

\begin{definicion}[Book-building]
El \textbf{book-building} es un proceso mediante el cual el banco de inversión determina un rango de precios y recopila indicaciones no vinculantes de interés de inversores, a menudo después de road-shows. El banco tiene discreción completa en la asignación de acciones.
\end{definicion}

\begin{definicion}[Subastas]
En las \textbf{subastas}, se establece un precio de equilibrio del mercado después de que se presenten las ofertas. El precio se determina mediante un mecanismo de mercado que iguala la oferta y la demanda.
\end{definicion}

\begin{proposicion}[Prevalencia del Book-building]
El 99\% de las OPIs en Estados Unidos y el 80\% de las OPIs a nivel mundial utilizan book-building.
\end{proposicion}

Esta proposición refleja la dominancia del método de book-building en los mercados de capitales modernos, particularmente en los mercados desarrollados.

\subsection{Ofertas de Precio Fijo}

\subsubsection{Características del Precio Fijo}

En las ofertas de precio fijo, el precio final se establece muy cerca de la fecha de la oferta real. Una vez establecido el precio, \textbf{no puede ser aumentado}, aunque puede ser reducido si las condiciones del mercado se deterioran.

\subsubsection{Determinación del Precio Inicial}

El precio inicial en ofertas de precio fijo se determina mediante:

\begin{enumerate}
\item \textbf{Análisis de flujos de caja descontados}: Valoración fundamental basada en proyecciones de flujos de caja futuros descontados a una tasa apropiada.
\item \textbf{Múltiplos de la industria}: Comparación con empresas similares utilizando ratios como:
   \begin{itemize}
   \item Precio/beneficios (\textit{price-to-earnings}, P/E).
   \item Mercado a valor contable (\textit{market-to-book}).
   \item Mercado a ventas (\textit{market-to-sales}).
   \end{itemize}
\end{enumerate}

Los múltiplos se utilizan ampliamente en OPIs debido a su simplicidad y a la disponibilidad de datos comparables de empresas públicas.

\subsubsection{Relación entre Tiempo y Rentabilidad del Primer Día}

\begin{proposicion}[Relación Tiempo-Rentabilidad en Precio Fijo]
Cuanto mayor es el tiempo que transcurre entre cuando se establece un precio de oferta fijo y cuando comienza la negociación, mayor es la rentabilidad promedio del primer día.
\end{proposicion}

\begin{proof}[Explicación de la Proposición]
Esta relación se debe a que cuanto mayor es el tiempo hasta la finalización, mayor es el riesgo de que las condiciones del mercado se deterioren y la oferta falle. Para reducir la probabilidad de una oferta fallida, se establece un precio de oferta más bajo. Condicional a que la oferta tenga éxito, el infravaloramiento esperado es relativamente alto.
\end{proof}

Esta proposición ilustra el trade-off entre riesgo y rentabilidad en el contexto de ofertas de precio fijo: mayor incertidumbre temporal requiere mayor infravaloramiento para compensar el riesgo de fallo de la oferta.

\subsection{Book-building}

El método de book-building es el más utilizado en los mercados desarrollados y permite una mayor flexibilidad en la determinación del precio y la asignación de acciones.

\subsubsection{Características del Book-building}

\begin{enumerate}
\item \textbf{Discreción completa del suscriptor}: Los suscriptores tienen discreción completa en la asignación de acciones, permitiéndoles favorecer a ciertos inversores (por ejemplo, inversores institucionales de largo plazo).
\item \textbf{Publicación de rango de precios}: Se publica un rango de precios (indicativo solamente). Los inversores no tienen que mantenerse dentro del rango, pero el rango proporciona una guía sobre el precio esperado.
\item \textbf{Campaña de marketing}: Como parte de la campaña de marketing, se realiza un road-show y educación de inversores.
\end{enumerate}

\subsubsection{Proceso de Recopilación de Indicaciones de Interés}

Durante el proceso de book-building, el banco de inversión obtiene indicaciones de interés de inversores institucionales sobre cuál es su demanda:

\begin{itemize}
\item Algunos inversores demandarán diferentes cantidades a diferentes precios.
\item Algunos establecerán límites de precio.
\item Estas indicaciones \textbf{no son vinculantes}, pero proporcionan información valiosa sobre la demanda del mercado.
\end{itemize}

\subsubsection{Asignación Histórica}

Históricamente, el criterio ha sido asignar una porción tanto a inversores institucionales como a inversores individuales (minoristas). Sin embargo, en muchos países, como Estados Unidos y el Reino Unido, la asignación está a discreción de los bancos de inversión, lo que permite buenos tratos para clientes preferidos.

\subsubsection{Establecimiento del Precio Final}

Un día antes de la oferta, el banco establece el precio final basado en el libro de órdenes. Cuando se establece el precio final, el banco de inversión llena las órdenes:

\begin{itemize}
\item Si la oferta está sobresuscrita, las asignaciones se reducen proporcionalmente.
\item En el proceso de book-building, la sobresuscripción es esencialmente determinada por los bancos de inversión; dado que tienen indicaciones de interés, conocen la demanda.
\end{itemize}

\subsection{Subastas}

Las subastas han sido utilizadas en muchos países, incluyendo Francia, Israel, Japón y Taiwán, para fijar precios y asignar OPIs.

\begin{proposicion}[Rentabilidad del Primer Día en Subastas]
Las subastas han estado asociadas con rentabilidades promedio del primer día bajas pero positivas. Generalmente, estas rentabilidades son menores que las de ofertas de precio fijo o book-building.
\end{proposicion}

\begin{ejemplo}[Subasta de Google en 2004]
Google realizó su OPI en 2004 utilizando una subasta holandesa (\textit{Dutch auction}). El precio de oferta fue de \$85 por acción. Este ejemplo ilustra cómo una empresa tecnológica importante utilizó un método alternativo de fijación de precios, aunque el método no se ha generalizado ampliamente después de este caso.
\end{ejemplo}

\section{Características Desconcertantes de las OPIs}

La investigación empírica sobre OPIs ha identificado tres características que desafían las explicaciones tradicionales de los mercados financieros eficientes. Estas características han generado extensa investigación académica y múltiples teorías explicativas.

\subsection{Tres Características Desconcertantes}

Las tres características empíricas principales que han desconcertado a los investigadores son:

\begin{enumerate}
\item \textbf{Infravaloramiento a Corto Plazo}: Las OPIs típicamente se venden por debajo de su valor de mercado en el primer día de negociación.
\item \textbf{Ciclos en el Número de OPIs}: Existen períodos de "mercados calientes" (\textit{hot issue markets}) y "mercados fríos" con variación significativa en el volumen de OPIs.
\item \textbf{Desempeño Inferior a Largo Plazo}: Las OPIs tienden a tener un desempeño inferior al de empresas comparables en los años posteriores a la emisión.
\end{enumerate}

\subsection{Infravaloramiento a Corto Plazo}

\subsubsection{Definición y Evidencia Empírica}

\begin{definicion}[Infravaloramiento del Primer Día]
El \textbf{infravaloramiento del primer día} se define como la situación en la cual el precio de cierre del mercado en el primer día de negociación de una OPI es mayor que el precio de oferta.
\end{definicion}

\begin{proposicion}[Universalidad del Infravaloramiento]
En todos los países con mercado de valores, las OPIs están infravaloradas. El infravaloramiento es un fenómeno universal observado en todos los mercados de capitales.
\end{proposicion}

\begin{definicion}[Dinero Dejado sobre la Mesa]
La cantidad de \textbf{dinero dejado sobre la mesa} (\textit{money left on the table}) se define como la diferencia entre el precio de cierre del primer día y el precio de oferta, multiplicada por el número de acciones vendidas:
\begin{equation}
\text{MLOT} = (P_{\text{cierre}} - P_{\text{oferta}}) \times N
\end{equation}
donde $P_{\text{cierre}}$ es el precio de cierre del primer día, $P_{\text{oferta}}$ es el precio de oferta, y $N$ es el número de acciones vendidas.
\end{definicion}

El dinero dejado sobre la mesa representa:
\begin{itemize}
\item La ganancia del primer día recibida por inversores a los que se les asignaron acciones al precio de oferta.
\item Una transferencia de riqueza de los accionistas de la empresa emisora a estos inversores.
\end{itemize}

\subsubsection{Evidencia Internacional}

La evidencia empírica muestra variación significativa en el grado de infravaloramiento entre países. Por ejemplo, países como China han mostrado infravaloramiento promedio extremadamente alto (aproximadamente 130-260\%), mientras que países como Dinamarca, Canadá, y los Países Bajos muestran infravaloramiento promedio más moderado (aproximadamente 5\%). Estados Unidos muestra un infravaloramiento promedio de aproximadamente 15\%.

\subsubsection{Objetivos en la Fijación de Precios}

Antes de explicar las razones del infravaloramiento, es importante notar que la minimización del infravaloramiento no es necesariamente óptima. Los objetivos potenciales incluyen:

\begin{enumerate}
\item \textbf{Maximizar el valor de la empresa}: Maximizar el precio del mercado secundario.
\item \textbf{Maximizar los ingresos para los accionistas vendedores}: Maximizar el precio de oferta.
\item \textbf{Minimizar los costes asociados con el proceso de salida a bolsa}: Minimizar el diferencial y el precio de oferta.
\end{enumerate}

Estos objetivos pueden estar en conflicto, y la empresa debe determinar qué objetivo priorizar.

\subsubsection{Explicaciones del Infravaloramiento}

Existen múltiples teorías que explican el infravaloramiento, y estas razones no son mutuamente excluyentes. Su importancia relativa difiere entre países, mecanismos contractuales y períodos de tiempo.

\paragraph{Maldición del Ganador (\textit{Winner's Curse})}

\begin{teorema}[Argumento de la Maldición del Ganador]
Un aumento en el precio es necesario para atraer a todos los inversores. De lo contrario, los inversores no informados podrían temer ser desfavorecidos, ya que los inversores informados adquirirían la mayoría de las acciones de empresas prometedoras, dejando a los inversores no informados con empresas no prometedoras.
\end{teorema}

\begin{proof}
Si los inversores no informados reciben todas las acciones que solicitan, es porque los inversores informados no quieren esas acciones. Esto implica que una fracción mayor de acciones sobrevaloradas se asignan a inversores no informados. Por tanto, los inversores no informados solo presentarán órdenes de compra si, en promedio, las OPIs están suficientemente infravaloradas para compensarlos por la asignación sesgada de nuevas emisiones. Si las OPIs se valoran de manera justa en promedio, los inversores no informados pierden. El mercado de OPIs requiere inversores no informados, de ahí la necesidad de infravaloramiento.
\end{proof}

\paragraph{Hipótesis de Señalización}

\begin{teorema}[Hipótesis de Señalización]
El infravaloramiento de OPIs permite a las empresas buenas señalizarse y obtener un mejor precio en ventas de capital posteriores. Los inversores saben que solo las empresas buenas pueden recuperar el coste de esta señal de emisiones posteriores. Es más costoso para las empresas malas imitar a las empresas buenas al bajar el precio de la OPI.
\end{teorema}

La racionalidad es evitar el impacto negativo de SEOs posteriores. Sin embargo, la evidencia empírica ha mostrado que la relación hipotética entre rentabilidades iniciales y éxito posterior en nuevas emisiones secundarias no está presente.

\paragraph{Corrupción}

Existe un problema de agencia entre los tomadores de decisiones en empresas emisoras y otros accionistas pre-emisión. Este problema contribuye a una disposición a contratar suscriptores que infravaloran la OPI. Un suscriptor que tiene otras "OPIs calientes" (altamente demandadas) puede hacer "pagos colaterales" a los tomadores de decisiones de una empresa emisora estableciendo cuentas de corretaje personales para estos individuos y luego asignando las OPIs calientes a estas cuentas.

\begin{ejemplo}[Caso Salomon Smith Barney]
En agosto de 2002, documentos fueron publicados por el Comité de Servicios Financieros de la Cámara de Representantes de Estados Unidos mostrando que Salomon Smith Barney (SSB) asignó OPIs calientes a directores ejecutivos de muchas empresas de telecomunicaciones durante 1996-2000. Durante este período, SSB tenía una gran participación de mercado en el negocio de suscripción de capital y fusiones y adquisiciones en esta industria.
\end{ejemplo}

\paragraph{Cascadas de Información}

\begin{definicion}[Cascada de Información]
Si los inversores potenciales prestan atención no solo a su propia información sobre una nueva emisión, sino también a si otros inversores están comprando, pueden desarrollarse \textbf{efectos de moda} o \textbf{cascadas de información}.
\end{definicion}

Si un inversor ve que nadie más quiere comprar, puede decidir no comprar incluso cuando posee información favorable. Para prevenir que esto suceda, un emisor puede querer infravalorar una emisión para inducir a los primeros inversores potenciales a comprar e inducir una cascada de todos los demás inversores. El objetivo es evitar cascadas de información negativas, de ahí el infravaloramiento.

\paragraph{Hipótesis de Evitación de Demandas}

La frecuencia y severidad de demandas futuras pueden reducirse mediante infravaloramiento, ya que solo los inversores que pierden dinero tienen derecho a daños. Sin embargo, infravalorar la OPI es una forma muy costosa de reducir la probabilidad de una demanda futura.

\begin{ejemplo}[OPI de Facebook en 2012]
La OPI de Facebook en 2012 resultó en más de 40 demandas presentadas respecto a la OPI en el mes que siguió. Un mes después de la OPI, el precio estaba alrededor de un 25\% por debajo del precio de oferta. Sin embargo, la evidencia empírica sugiere que no hay diferencia en infravaloramiento entre OPIs seguidas de demandas y aquellas que no lo fueron.
\end{ejemplo}

\paragraph{Otras Explicaciones}

\begin{itemize}
\item \textbf{Hipótesis de Propiedad Dispersa}: Los directivos no les gustan los grandes accionistas mayoritarios (amenaza). El infravaloramiento implica sobresuscripción. Racionamiento a favor de pequeños accionistas.
\item \textbf{No hay infravaloramiento (comportamental)}: Lo que se observa es una burbuja especulativa. El precio de oferta es un precio justo; el primer precio del mercado secundario es demasiado alto. Implicación: el precio debería disminuir cuando la burbuja estalla.
\end{itemize}

\subsubsection{Evidencia Empírica del Infravaloramiento}

\begin{proposicion}[Patrones Empíricos del Infravaloramiento]
\begin{enumerate}
\item Se observa mayor infravaloramiento para empresas sobre las cuales hay mayor incertidumbre, es decir, empresas más pequeñas y empresas que son más volátiles en el mercado secundario.
\item Las empresas (y sus suscriptores) infravaloran la emisión incluso cuando saben que será masivamente sobresuscrita.
\item La opción Green Shoe (o de sobredemanda) reduce la Maldición del Ganador al proporcionar un mecanismo para estabilizar el precio.
\end{enumerate}
\end{proposicion}

\subsection{Ciclos en el Número de OPIs}

\subsubsection{Evidencia de Ondas de OPIs}

Existen períodos de "mercados calientes" (\textit{hot issue periods}) y períodos de "mercados fríos" (\textit{cold issue periods}). Los mercados calientes pueden no ser sorprendentes: las empresas acuden al mercado cuando necesitan financiación, y esas necesidades pueden seguir el ciclo económico.

\begin{proposicion}[Variación en el Número de OPIs]
Sin embargo, la variación parece demasiado grande para ser explicada por el ciclo económico. Las variables del ciclo económico explican poco de la variación en el número de OPIs.
\end{proposicion}

La evidencia empírica muestra variación significativa en el número de OPIs a lo largo del tiempo. Por ejemplo, en Estados Unidos durante 1980-2000, un promedio de 310 empresas salieron a bolsa cada año, mientras que en 2001-2012, el promedio fue de solo 99 empresas por año.

\subsubsection{Teorías Explicativas}

\paragraph{Teoría \#1: La Historia del Mercado Irracional}

\begin{proposicion}[Teoría del Mercado Irracional]
Esta teoría sostiene que el comportamiento del mercado está impulsado por emociones humanas:
\begin{itemize}
\item Durante períodos optimistas, las personas comprarán cualquier cosa.
\item Durante períodos pesimistas, las personas comprarán muy poco.
\end{itemize}
Los mercados fluctúan entre estados "calientes" y "fríos" basados en la "*sentiment* del inversor". Esta teoría no es ampliamente aceptada por economistas de "mercados eficientes".
\end{proposicion}

\paragraph{Teoría \#2: La Historia del Mercado Eficiente}

\begin{teorema}[Teoría del Mercado Eficiente]
En ciertos momentos hay grandes oportunidades de inversión. Eso es cuando el mercado de OPIs está "caliente"; invertir en estas empresas es excelente ya que van a usar el dinero bien. Por el contrario, en otros momentos no hay muy buenas oportunidades de inversión. Los inversores creerán racionalmente que una empresa que intenta emitir en tal momento debe estar desesperada y, por tanto, no pagarán mucho por las acciones de la empresa. Eso es cuando el mercado está frío. Bajo esta teoría, las empresas tienden a emitir acciones al público en momentos en que hay grandes oportunidades de inversión, ya que es cuando los inversores están dispuestos a comprar y pagar más.
\end{teorema}

\subsection{Desempeño Inferior a Largo Plazo}

\subsubsection{Evidencia Empírica}

Las rentabilidades anuales de las OPIs en los 5 años posteriores a la emisión se reportan junto con un benchmark utilizando emparejamiento por tamaño o estilo. Los emparejamientos se basan en tamaño (capitalización de mercado) y estilo (valor contable a mercado).

\begin{proposicion}[Desempeño Inferior de las OPIs]
Las OPIs han tenido un desempeño inferior a otras empresas del mismo tamaño en un promedio del 5.2\% por año durante los 5 años posteriores a la emisión, sin incluir la rentabilidad del primer día.
\end{proposicion}

Esto parece ir en contra de mercados eficientes y racionales. Sin embargo:
\begin{itemize}
\item Todavía no sabemos si este es un fenómeno generalizado.
\item No tenemos una teoría consistente para explicar este patrón sistemático.
\end{itemize}

\subsubsection{Hipótesis Explicativas (Sabor Comportamental)}

\paragraph{Hipótesis de Divergencia de Opinión (Miller 1977)}

\begin{teorema}[Hipótesis de Divergencia de Opinión]
Los inversores más optimistas compran OPIs. Si hay alta incertidumbre sobre el valor de una OPI, las valoraciones optimistas serán mucho más altas que las pesimistas. A medida que se dispone de más información, esta divergencia se estrecha, causando que los precios caigan y las OPIs tengan un desempeño inferior a largo plazo.
\end{teorema}

\paragraph{Hipótesis del Impresario (Shiller 1990)}

\begin{teorema}[Hipótesis del Impresario]
Los banqueros de inversión (impresarios) infravaloran las OPIs para crear una apariencia de exceso de demanda, similar a un promotor de conciertos de rock creando un "evento". Shiller predice que las empresas con las mayores rentabilidades iniciales tendrán posteriormente las menores rentabilidades.
\end{teorema}

\paragraph{Hipótesis de Ventanas de Oportunidad}

\begin{teorema}[Hipótesis de Ventanas de Oportunidad]
Durante períodos en que los inversores son particularmente optimistas sobre el potencial de crecimiento de las empresas que salen a bolsa, los grandes ciclos en volumen representan empresas "cronometrando" sus OPIs para aprovechar el sentimiento del inversor. Predice que las empresas que salen a bolsa en períodos de alto volumen tienen más probabilidades de estar sobrevaloradas que otras OPIs, y estos períodos de alto volumen deberían estar asociados con las menores rentabilidades a largo plazo. Este patrón efectivamente existe.
\end{teorema}

\section{Síntesis y Conclusiones}

Este documento ha presentado un tratamiento riguroso de las Ofertas Públicas Iniciales, desde sus definiciones fundamentales hasta las características empíricas que desafían las explicaciones tradicionales de los mercados financieros.

Las OPIs representan un mecanismo crucial mediante el cual las empresas privadas acceden a los mercados públicos de capital, transformando su estructura de propiedad y financiación. La decisión de realizar una OPI implica una evaluación cuidadosa de múltiples factores, incluyendo beneficios de acceso a capital y liquidez versus costes de cumplimiento regulatorio y pérdida de control.

El proceso de OPI es complejo y requiere la coordinación de múltiples participantes, con los bancos de inversión desempeñando un papel central en la valoración, fijación de precios, y comercialización de las acciones. El método de book-building domina los mercados desarrollados, proporcionando flexibilidad en la determinación del precio y asignación de acciones.

Las tres características desconcertantes de las OPIs—infravaloramiento a corto plazo, ciclos en el número de OPIs, y desempeño inferior a largo plazo—continúan generando investigación académica y debate teórico. Múltiples explicaciones han sido propuestas, desde teorías basadas en información asimétrica hasta explicaciones comportamentales, reflejando la complejidad de los mercados financieros y la interacción entre participantes racionales y factores psicológicos.

La comprensión de estos fenómenos es esencial para inversores, empresas emisoras, y reguladores que buscan optimizar el funcionamiento de los mercados de capitales y maximizar el valor para todas las partes involucradas.

\end{document}

